\documentclass[10pt]{amsart}
\usepackage[a4paper,margin=23mm]{geometry}
\usepackage{amsmath,amssymb,amsthm,mathtools,hyperref,graphicx,xspace,letltxmacro,xpatch}
\usepackage[english]{babel}
\newtheorem{theo}{Theorem}[section]
\newtheorem{lemm}[theo]{Lemma}\newtheorem{prop}[theo]{Proposition}\newtheorem{coro}[theo]{Corollary}\newtheorem{conj}[theo]{Conjecture}
\theoremstyle{definition}\newtheorem{defi}[theo]{Definition}\newtheorem{rema}[theo]{Remark}\newtheorem{exam}[theo]{Example}\newtheorem{exem}[theo]{Example}\newtheorem{nota}[theo]{Notation}
\newcommand{\ie}{i.e.}\newcommand{\eg}{e.g.}\newcommand{\resp}{resp.}\newcommand{\skpt}{\par\smallskip\noindent}
\emergencystretch=3em\pagestyle{plain}

\usepackage{fancybox}
\usepackage{framed}
%\usepackage{enumerate}



 	


\newcommand{\RR}{\mathbb{R}}
\newcommand{\NN}{\mathbb{N}}
\newcommand{\CC}{\mathbb{C}}
\newcommand{\TT}{\mathbb{T}}
\renewcommand{\SS}{\mathbb{S}}

\newcommand{\Cc}{\mathcal{C}}
\newcommand{\Ac}{\mathcal{A}}
\newcommand{\Bc}{\mathcal{B}}
\newcommand{\Fc}{\mathcal{F}}
\newcommand{\Lc}{\mathcal{L}}
\newcommand{\Pc}{\mathcal{P}}
\newcommand{\Qc}{\mathcal{Q}}
\newcommand{\Oc}{\mathcal{O}}
\newcommand{\drm}{\mathrm{d}}

\DeclarePairedDelimiter{\VERT}{|\mkern -1.6mu|\mkern -1.6mu|}{|\mkern -1.6mu|\mkern -1.6mu|}

%\newcommand{\Kc}{\mathcal{K}}
%\newcommand{\Uc}{\mathcal{U}}
%\newcommand{\Vc}{\mathcal{V}}
%\newcommand{\Wc}{\mathcal{W}}
%\newcommand{\Nc}{\mathcal{N}}


\newcommand{\grad}{\nabla}
\renewcommand{\div}{\operatorname{div}}
\newcommand{\cnvh}{convhull}

%{c}
\newcommand{\vect}[2]{\left(
\begin{matrix} #1 \\ #2
\end{matrix}\right)}
\newcommand{\trans}{{^{\intercal}}}

\newcounter{stepcount}
%%\renewcommand{\thestepcount}{\alph{stepcount}}
\newenvironment{step}{\refstepcounter{stepcount}\begin{proof}[Step~\thestepcount]}{\let\qed\relax\end{proof}}



%\newdimen\texpscorrection
%\texpscorrection=0truecm
%\newdimen\figcenter
%\newcommand*\figurewithtex [1] #2 #3 #4 #5\\\null \goodbreak\figcenter=\hsize\relax
%\advance\figcenter by -#4truecm
%\divide\figcenter by 2
%
%\begin{figure}[hbt]
%\vskip #3truecm\noindent\hskip\figcenter
%\special{psfile=#1}{\hskip\texpscorrection\failedinput{\baselineskip=0.8\baselineskip
%\noindent \vbox{\noindent {\footnotesize #5}}\end{figure}}}
%\newcommand*\point[1] #2 #3 {\rlap{\kern #1 truecm\raise #2 truecm \hbox{#3}}}




%%%%%%%%%%%%%%%%%%%%%%%%%%%%%%%%%%%%%%%%%%%%%%%%%


\graphicspath{{./figures/}}

\newcommand*{\mk}{\mkern -1mu}
\newcommand*{\Mk}{\mkern -2mu}
\newcommand*{\mK}{\mkern 1mu}
\newcommand*{\MK}{\mkern 2mu}

\hypersetup{urlcolor=purple, linkcolor=blue, citecolor=red}

\newcommand*{\romanenumi}{\renewcommand*{\theenumi}{\roman{enumi}}}
\newcommand*{\Romanenumi}{\renewcommand*{\theenumi}{\Roman{enumi}}}
\newcommand*{\alphenumi}{\renewcommand*{\theenumi}{\alph{enumi}}}
\newcommand*{\Alphenumi}{\renewcommand*{\theenumi}{\Alph{enumi}}}
\let\oldtilde\tilde
\renewcommand*{\tilde}[1]{\mathchoice{\widetilde{#1}}{\widetilde{#1}}{\oldtilde{#1}}{\oldtilde{#1}}}\let\oldexists\exists
\renewcommand*{\exists}{\mathrel{\oldexists}}
\let\oldforall\forall
\renewcommand*{\forall}{\mathrel{\oldforall}}

%%%%%%%%%%%%%%%%%%%%%%%%%%%%%%%%%%%%%%%%%%%%%%%%
\newtheorem{lemma}[theo]{Lemma}
\newenvironment{remas}[1][]{\begin{rema}[#1]}{\end{rema}}
